\section{Introduction}
% =====================================================================

Spatially resolved transcriptomics reports gene expression together with the
physical coordinates of each measurement \citep{stahl2016visium,chen2015merfish}.
Nearly all machine learning built on it---spatial domain identification,
clustering, imputation, denoising---represents a tissue section as a graph over
the measured locations and applies message passing
\citep{hu2021spagcn,dong2022stagate,palla2022squidpy}. This matches how the data
arrive, but it fixes a strong inductive bias: the tissue \emph{is} the sampling
lattice. Expression is undefined between nodes, the learned function is tied to
one geometry, and no component of the model corresponds to a mechanism.

Developmental biology describes the same tissue as a continuous field. Morphogens
are produced, diffuse, decay and interact, and cells read the resulting
concentration landscape \citep{turing1952,wolpert1969positional,gierer1972theory}.
Reaction--diffusion models of this form account for phenomena from digit
patterning \citep{sheth2012hox} to Nodal/Lefty signaling \citep{muller2013nodal},
with measured gradient lengths of tens to hundreds of microns
\citep{kicheva2007kinetics,yu2009fgf8}. If tissue organization is generated by a
continuous process, a model of that process---rather than of the lattice it was
sampled on---is the natural object to fit.

We introduce \textbf{Neural Morphogen Operators} (NMO), which encode irregularly
sampled expression into a continuous latent field, evolve that field under a
learned anisotropic reaction--diffusion PDE, and decode expression at arbitrary
coordinates (Fig.~\ref{fig:overview}). The formulation is designed so that the
learned dynamics are both numerically well behaved and physically readable: the
diffusion half-step is solved exactly in the Fourier domain, which removes any
step-size restriction and renders the learned tensor $\mD_\theta$ directly
interpretable as a diffusion length in microns.

\paragraph{Contributions.}
\begin{enumerate}\itemsep1pt \parskip0pt
\item We introduce a PDE-constrained neural operator for spatial omics, coupling
  a graph kernel-integral encoder, a spectral reaction--diffusion operator and a
  coordinate-free continuous decoder (Section~\ref{sec:method}).
\item We show that the resulting integrator is unconditionally stable and
  second-order accurate, with structurally guaranteed positive-definite
  diffusion and a uniform absorbing set for the field
  (Propositions~\ref{prop:contraction}--\ref{prop:bounded}).
\item We show that the conventional physics-informed residual is uninformative
  for an exactly integrated operator, and derive the steady-state constraint that
  replaces it (Proposition~\ref{prop:vacuity}).
\item We demonstrate improved masked-region reconstruction over graph and
  Gaussian-process baselines across four spatial technologies, and isolate the
  operator as the source of the gain using an exactly parameter-matched control.
\item We establish that a single stationary snapshot determines the reaction
  network only on a Lebesgue-null subset of latent space
  (Theorem~\ref{thm:identifiability}), which delimits what any method of this
  class can identify and predicts the outcome of our perturbation experiment.
\end{enumerate}

% =====================================================================
